\subsection{Aplicación de salida y condicionamiento del registro}
\label{vii:sec:salida-condicionamiento}

La inscripción anterior se compone con el lector del aparato sobre el
estado genealógico completo. El acoplamiento conserva la operación que
produce la correlación; la aplicación de salida expresa el resultado; el
condicionamiento restringe las prolongaciones compatibles con ese registro.

Sean \((\Omega_S,\Sigma_S)\) y \((\Omega_M,\Sigma_M)\) los espacios
medibles de estados genealógicos del sistema y del aparato, obtenidos de
sus cilindros y registros compatibles. Un contexto de medición se escribe
\(\mathfrak M_a=(U_a,q_a,\mathcal C_a,m_0)\): \(m_0\) es la
preparación del aparato, \(\mathcal C_a\) fija las compatibilidades y el
horizonte de prolongación, y
\[
 U_a:\Omega_S\times\Omega_M\longrightarrow\Omega_S\times\Omega_M,
 \qquad q_a:\Omega_M\longrightarrow Y_a
\]
son, respectivamente, el acoplamiento medible y el lector medible hacia
el espacio de resultados \((Y_a,\Sigma_{Y_a})\). En un régimen reversible,
\(U_a\) es biyectivo y bimedible; en la realización de Hilbert se exige
asimismo la unitariedad correspondiente. La salida individual es la composición
\begin{equation}
 \operatorname{Out}_a(x)
 =q_a\bigl(\operatorname{pr}_M U_a(x,m_0)\bigr).
 \label{vii:eq:out-genealogico}
\end{equation}
Queda determinada por el estado completo, la preparación del aparato y
el contexto. La descripción estadística de una preparación constituye
otra lectura del conjunto de estados compatibles.

\begin{proposition}[Medida de resultados y restricción a la fibra]
\label{vii:prop:medida-salida}
Sea \(\Omega_a\in\Sigma_S\) el conjunto de estados compatibles con la
preparación y el contexto, provisto de una probabilidad \(\mu_a\) sobre
\((\Omega_a,\Sigma_S|_{\Omega_a})\).
La aplicación \(\operatorname{Out}_a|_{\Omega_a}\) es medible y produce
la probabilidad de resultados
\begin{equation}
 \nu_a(B)=\mu_a\bigl(\operatorname{Out}_a^{-1}(B)\cap\Omega_a\bigr),
 \qquad B\in\Sigma_{Y_a}.
 \label{vii:eq:medida-imagen-salida}
\end{equation}
Para un resultado \(y\) con singleton medible, la fibra
\(\Omega_{a,y}=\{x\in\Omega_a:\operatorname{Out}_a(x)=y\}\)
es medible. Cuando \(\mu_a(\Omega_{a,y})>0\),
\begin{equation}
 \mu_{a,y}(E)=
 \frac{\mu_a(E\cap\Omega_{a,y})}{\mu_a(\Omega_{a,y})}
 \label{vii:eq:condicionamiento-salida}
\end{equation}
para \(E\in\Sigma_S|_{\Omega_a}\) define una probabilidad concentrada
en esa fibra.
\end{proposition}
\begin{proof}
La inclusión \(x\mapsto(x,m_0)\), el acoplamiento, la proyección sobre
el aparato y su lector son medibles; su composición también lo es.
Las preimágenes conservan uniones disjuntas y el conjunto total, por lo
que \eqref{vii:eq:medida-imagen-salida} es numerablemente aditiva y tiene
masa uno. La fibra es la preimagen del singleton \(\{y\}\) dentro de
\(\Omega_a\). La restricción de \(\mu_a\) a esa fibra es una medida
positiva y numerablemente aditiva; dividir por su masa positiva la normaliza.
\end{proof}

Para un prefijo \(s\) de profundidad \(n\), las prolongaciones posteriores
al registro forman
\[
 \operatorname{Fut}_{a,y}(s)
 =\{x\in\Omega_{a,y}:x|_n=s\}.
\]
La restricción conserva el prefijo ya recorrido y el registro del
acoplamiento; modifica la colección de prolongaciones admitidas. La salida
\eqref{vii:eq:out-genealogico} actúa sobre estados completos; la ley de
resultados \eqref{vii:eq:medida-imagen-salida} se obtiene transportando
la medida de preparación por esa aplicación. La representación por una
matriz de densidad expresa el nivel
estadístico de la realización cuántica y conserva sus condiciones de
correspondencia con las fibras genealógicas.
